\section{Background}

\subsection{ZNS Placement and Write Amplification}

\zns SSDs expose a storage interface in which each zone is written sequentially and later reset as a unit~\cite{zns}.  This interface removes part of the hidden device FTL policy and gives the host more control over placement.  The benefit depends on grouping objects with compatible lifetimes.  When a zone contains only expired objects, reset is cheap.  When a zone also contains live objects, the host must copy those valid blocks elsewhere before reset, increasing write amplification factor (WAF).

Prior placement systems therefore focus on predicting block invalidation time.  \sepbit separates data using historical invalidation behavior.  \midas adapts group configuration to workload patterns.  \dogi combines hot filtering, learned invalidation-time prediction, GC relocation policy, and dynamic group organization.  The DOGI paper evaluates both static-skew workloads such as FIO Zipf and YCSB-A/F and dynamic workloads such as Varmail, Alibaba, and Exchange~\cite{dogi}.  This is the right baseline family for \system because these systems represent the strongest storage-history approach.

\subsection{Post-Quantum Object Lifetimes}

Post-quantum workloads introduce a different kind of lifetime signal.  A KEM artifact or ephemeral secret is not short-lived because it has a particular LBA pattern; it is short-lived because the cryptographic protocol invalidates it.  A signature log may be append-only and long-lived, while a KMS envelope may die when an epoch closes.  A certificate metadata object may follow a rotation period rather than a session boundary.  These differences exist within the same service and often within the same burst of I/O.

\begin{table}[t]
\centering
\caption{Post-quantum objects have different lifetime and security behavior even when they are generated by the same service.}
\label{tab:pqc-objects}
\footnotesize
\begin{tabularx}{\columnwidth}{lYY}
\toprule
\textbf{Object} & \textbf{Lifetime driver} & \textbf{Placement implication} \\
\midrule
KEM artifact & Session or epoch close & Group by epoch; reset quickly. \\
Ephemeral secret & Session close & Isolate from payload and logs. \\
KMS envelope & Key epoch & Pack with same epoch family. \\
Certificate metadata & Rotation policy & Use rotation-zone family. \\
Signature log & Retention policy & Append-only placement. \\
Payload & Application update pattern & Use history-based fallback. \\
\bottomrule
\end{tabularx}
\end{table}

Table~\ref{tab:pqc-objects} gives the placement consequence.  A single hotness label cannot safely place these objects because two objects with similar write frequency may have opposite reclaim requirements.  Conversely, two objects with different LBAs and different sizes may be ideal zone neighbors if they die at the same epoch boundary.

\subsection{Claim Boundary}

\system is not an oracle.  It does not receive a future deletion timestamp for every object, and it does not assume that every hint is perfect.  The hint expresses lifecycle class and epoch membership, not secret key material.  Correctness is protected by a conservative epoch manager: a zone family is reset only after the manager confirms that all objects in that family have expired or have been safely migrated.

The security claim is also scoped.  Zone reset makes a zone reusable and can reduce exposure, but physical sanitization is device-dependent.  Our physical device advertises crypto-erase and block-erase sanitize support, and the crypto-erase sanitize command path has been executed successfully.  Nevertheless, the default claim is reset eligibility and exposure-window reduction.  A physical erase claim requires issuing sanitize or equivalent hardware crypto-erase semantics at epoch boundaries.
